\documentclass[11pt]{article}
\usepackage[margin=1in]{geometry}
\usepackage{amsmath,amssymb,amsthm,mathtools}
\usepackage{hyperref}
\setlength{\emergencystretch}{2em}
\newtheorem{theorem}{Theorem}
\newtheorem{lemma}[theorem]{Lemma}
\newtheorem{proposition}[theorem]{Proposition}
\newcommand{\Q}{\mathbb Q}
\newcommand{\Z}{\mathbb Z}
\newcommand{\R}{\mathbb R}
\newcommand{\C}{\mathbb C}
\newcommand{\pr}{\operatorname{pr}}
\newcommand{\Tr}{\operatorname{Tr}}
\newcommand{\Spin}{\operatorname{Spin}}
\newcommand{\SU}{\operatorname{SU}}
\title{The actual M45 integral Seifert congruence\\
An independent proof with compact exact checks}
\author{Campaign 9000059: Engineer submission}
\date{November 2026}
\begin{document}
\maketitle

\paragraph{Result and certification boundary.}
This submission proves the integral-congruence proposition in Section 9
of the staged four/five manuscript, using its \emph{actual} lattices.
The previously independently supported claims \texttt{C059BASE},
\texttt{C059APPLICATION}, and \texttt{M45:thm:tube} are inputs;
their local-algebra and tube arguments are not reopened.
The new work supplies the plane isometry, two-sided transfer adjunction,
the required semisimplicity, and the complete arithmetic passage.
In particular, it does not replace the integral tube lattice by its
rational orthogonal sum.

The proof is mathematical prose, with exact supporting calculations in
\texttt{verify.py}. It is neither a Lean proof nor a computation of a
huge integral Seifert matrix. Strong approximation gives an existence
proof, not an explicit matrix. Independent Reviewer acceptance and
independent attribution of the new literature inputs remain required.
Neither ambient equivalence, the separate two/three construction,
verified novelty, nor campaign completion is asserted here.

\section{Precise statement and retained inputs}

Fix an integer \(r\ge4\) and an integer \(n_0\) such that
\[
 n_0>5\cdot2^r,\qquad n_0\equiv1\pmod{2^r},\qquad n_0\equiv2\pmod5.
\]
Put
\[
 n_1=n_0+5\cdot2^r,\quad
 a=(n_0^2-1)/4,\quad b=(n_1^2-1)/4,\quad
 g=\gcd(a,b),\quad m=4ab+a+b.
\]
These choices satisfy the manuscript's strict tube inequalities.
Indeed \(n_0<n_1<2n_0\) gives \(a<b<4a\), so comparability is
uniform, and \(a>12\). For the manuscript's
\(\alpha=4b/(a+b)\), \(\beta=4a/(a+b)\), one has
\(\alpha+\beta=4\), \(a\alpha=b\beta\), and
\((ab-2b)/(a+b)>a/2-2>4\), with the other inequality identical
after interchanging the subtracted variable. Choose the common positive
\(s\) sufficiently small for the reviewed tube identification.
The Weierstrass units in the actual four-variable germs are retained.

Here are the plane maps, with no change of their projection data:
\[
\begin{split}
 h(X,T)&=T(X-T^2)(X-2T^2),\\
 q_1(u,w)&=(u-w-w^8)(u-w-2w^8)(u-w^2),\quad q_2=q_1-u^2,\\
 \pi_0(u,w)&=(u^2+q_1(u,w),u),\\
 \pi_1(x,w)&=(x,u(u-w^2)q_2(u,w)),\qquad u=x+w,\\
 d_\nu&=h\circ\pi_\nu.
\end{split}                                                    \tag{1}
\]
Let \(H=H_1(F_h;\Z)\), \(D_\nu=H_1(F_{d_\nu};\Z)\),
\(P_\nu=\pi_{\nu*}\), \(T_\nu=\pi_\nu^!\), and \(K_\nu=\ker P_\nu\).
The accepted base calculation gives degree four for both maps, rank
six for \(H\), and rank 51 for each \(D_\nu\).
All Seifert forms below have the same ordered positive-neighbor
filling convention as the reviewed tube result.

For the cone \(yz=w^4\), put \(P=w(y^a+z^b)\),
\(E=H_1(F_P;\Z)\), and let \(\tau:E\to\Z/4\) be the positive
\(y\)-angle class. The rational pairing \(e\) is one fourth the
upstairs Seifert pairing of full transfers. The accepted tube theorem
identifies the actual lattices with
\[
\begin{split}
 F_\nu&=K_{\nu,\Q}\oplus(H\otimes E)_\Q,\\
 s_\nu&=(-s_{d_\nu}|K_\nu)\perp(s_h\otimes e),\\
 \Lambda_\nu&=\{(\pr_\nu l,z):l\in D_\nu,\ z\in H\otimes E,\
        P_\nu l\equiv(1\otimes\tau)z\pmod{4H}\},\\
 \pr_\nu&=1-\tfrac14 T_\nu P_\nu. 
\end{split}                                                    \tag{2}
\]

\begin{theorem}\label{thm:result}
For every parameter choice just specified and both actual prepared
germs to which the reviewed tube theorem applies,
\((\Lambda_0,s_0)\) and \((\Lambda_1,s_1)\) are integrally isometric.
No compatibility of the plane isometry with \(P_0,P_1\) is assumed.
\end{theorem}

\section{The plane comparison, without identifying the projections}

\begin{lemma}\label{lem:plane}
There is an integral ordered Seifert isometry
\(I:(D_0,s_{d_0})\to(D_1,s_{d_1})\).
On \(H_\Q\) the monodromy is the sum of two trivial representations
and one copy of the four-dimensional fifth cyclotomic representation.
On either \(D_{\nu,\C}\), monodromy is semisimple at \(1\) and at
the primitive fifth roots.
\end{lemma}

\begin{proof}
In the first graph use \(t=u,\ x=u-w\); in the second use \(t=x+w\).
The six branches in either case have the following descriptions:
\[
 \begin{aligned}
 x&=A_e(t)=e t^8+O(t^9)\quad(e=1,2),\\
 t&=B_f(x)=f x^2+O(x^3)
 \quad\left(f=0,1,\frac{1+\sqrt{-3}}2,\frac{1-\sqrt{-3}}2\right).
 \end{aligned}
                                                               \tag{3}
\]
For the first graph the two \(A\)-remainders actually start at order
15. To check the only less immediate branches, substitute
\(u=v w^2\) in \(q_2\):
\[
 q_2(vw^2,w)=-(v^2-v+1)w^4+O(w^5),\qquad
 q_2(u,0)=u^2(u-1).
\]
The two simple roots give exactly the two remaining local branches.
In the second graph \(\psi(0,w)=-w^4+O(w^5)\); the factors
\(x-e\psi^2\) have a unique smooth branch by the implicit function
theorem and give the two \(A\)-branches in (3).
The other four branches are those of \(u(u-w^2)q_2\).
Thus both functions are a unit times the product in (3).

Here is a concrete justification that discarding the higher terms in
(3) preserves the \emph{integral ordered} form. Interpolate each
remainder to zero and interpolate the units through nonvanishing
units. Compactness of the interpolation interval makes all the
displayed order estimates uniform. Blow up the origin. The two
different tangent directions separate the two clusters. In the first
cluster the successive charts are \(x=t^k \xi_k\), \(1\le k\le8\).
Until the last chart the two strict transforms have the same section;
in the last chart they meet the exceptional divisor at \(1\) and \(2\).
In the other cluster the charts \(t=x^k\eta_k\), \(1\le k\le2\),
end at the four distinct values of \(f\) in (3).
All centres, transverse intersections, and component multiplicities
therefore persist throughout the interpolation. This is a simultaneous
normal-crossing resolution, not merely equality of Milnor numbers.

For clarity, the topological consequence used here can be seen in
these charts. Take small normal disk bundles around the exceptional
components, joined by bidisks at their crossings. Near a crossing
the total transform is \(u_\theta z_1^c z_2^d\), where \(u_\theta\)
is a unit and the positive integers \(c,d\) are independent of the
interpolation parameter \(\theta\). Away from crossings it is
\(u_\theta z_1^c\). Trivialize the moving transverse sections in
the disk bundles. The change in the argument of a unit on a small
chart has a continuous logarithm, so an angular coordinate corrects
it without changing the positive meridian degree. On overlaps the
correction fields are tangent to the normal-crossing strata; a
partition of unity combines fields satisfying the same linear
equation for the phase derivative. Their flows give the
value-angle fibration on the boundary of this plumbed neighborhood,
with positive binding meridians, throughout \(\theta\).
The transverse branch annuli between this boundary and an ordinary
Milnor boundary give collars: a level-tangent field extends an
outward field on each annulus, and away from these annuli the
absolute value of the function has a positive minimum.
Shrinking the value disk removes the remaining compact region.
Consequently these boundary fibrations are the ordinary Milnor
open books. The resulting oriented identification preserves linking
with a positive neighboring page, hence the integral ordered
Seifert form. It proves the assertion about \(I\) without saying
anything about its relation to the two projections.

The order-five action on the base page is
\((X,T)\mapsto(\zeta^2X,\zeta T)\). Every nonidentity element acts
freely. A cell structure lifted from the quotient has chain trace
zero for each such element. The alternating trace on homology
equals the chain trace, \(H_0\) has trace one, and \(H_2=0\).
Thus the trace on \(H_1\) is one. Rank six, together with the
rational irreducibles of \(C_5\), gives
\(u+4v=6,\ u-v=1\), or \(u=2,\ v=1\).

It remains to justify the Jordan assertion rather than infer it
from an eigenvalue count. Use the leading product
\[
 d_*=(x-t^8)(x-2t^8)t(t-x^2)(t^2-tx^2+x^4).
                                                               \tag{4}
\]
On the \(t\)-side its phase has weighted model
\(t^4(x-t^8)(x-2t^8)\), and on the \(x\)-side its weighted model
is \(x^2t(t-x^2)(t^2-tx^2+x^4)\).
The respective angular velocities, in turns per phase circuit, are
\((8,1)/20\) and \((1,2)/10\).
The correction factors are \(1+O(|x|^2/|t|)\) in
\(|x|\le2|t|\) and \(1+O(|t|^8/|x|)\) in \(|t|\le2|x|\);
their angular derivatives are uniformly small.
Logarithms and radial-ratio cutoffs straighten the phase in the
overlap to \(2\theta_x+4\theta_t\).
Both velocity fields differentiate this phase at rate one.
Their angular convex interpolation is therefore still transverse;
normalizing a small perturbation preserves this strict positivity.
Near the binding use the appropriate weighted field, which is
tangent to its branch. This verifies that the straightening is an
isotopy of open books, not just of their binding links.

The level \(2\theta_x+4\theta_t=\mathrm{constant}\) has two annuli.
Orient both cores \(\gamma_1,\gamma_2\) by increasing \(\theta_t\);
one circuit exchanges them without reversing these orientations.
After twenty circuits the rotations at the two ends are
\((8,1)\) and \((2,4)\). Their difference is
\((-6,3)=3(-2,1)\), three times the primitive phase-kernel circle.
Thus the twentieth return is the same signed threefold twist on
each joining annulus, modulo twists around actual page boundaries.
This statement follows directly by following the free angle across
each product annulus; the two annuli have identical oriented
coordinates. Boundary twists act trivially on first homology.

The oriented boundary of the \(t\)-side subsurface shows that
\(\gamma_1+\gamma_2\) is a sum of page-boundary classes and hence
is in the radical of intersection. The elementary crossing formula
for an annular twist is \(x\mapsto x+k(x\cdot\gamma)\gamma\):
each crossing acquires \(k\) copies of the annulus core.
Since the two \(k\)'s agree, it follows that, for the monodromy \(M\),
\[
 \operatorname{im}(M^{20}-1)\subset
   \Q(\gamma_1-\gamma_2),\qquad
 M(\gamma_1-\gamma_2)=-(\gamma_1-\gamma_2).
                                                               \tag{5}
\]
In particular \((M+1)(M^{20}-1)=0\).
The polynomial \((t+1)(t^{20}-1)\) has a repeated root only at
\(-1\). It has simple roots at \(1\) and all primitive fifth roots,
which proves the required semisimplicity.
\end{proof}

\begin{lemma}\label{lem:adjunction}
For the actual finite maps,
\[
 P_\nu T_\nu=4,\qquad
 s_{d_\nu}(T_\nu a,l)=s_h(a,P_\nu l),\qquad
 s_{d_\nu}(l,T_\nu a)=s_h(P_\nu l,a).              \tag{6}
\]
Hence \(D_{\nu,\Q}=K_{\nu,\Q}\perp T_\nu H_\Q\), in both
orders, and the projection is exactly the one in (2).
\end{lemma}

\begin{proof}
Use finite proper graph representatives over a sufficiently small
base cut. These cuts need not initially be balls. Put an ordinary
Milnor ball inside each. The six branch parametrizations above
give disjoint transverse annuli in the intervening closed region.
An outward annular field extends, by the normal complex
submersion, to a field tangent to nearby function levels.
Outside the resulting collars the function has a positive absolute
minimum. A smaller value disk therefore makes the larger tube an
ordinary acyclic tube with collars attached. This supplies acyclic
totals for the filling definition; acyclicity is not inferred just
from finiteness of the map.

A base cycle can be represented by loops avoiding the finitely
many branch values on a page. Their full preimage defines \(T_\nu\).
When a loop crosses a branch value, the full preimage of a small
disk supplies the bounding chain, so the construction depends only
on its homology class. Equivalently, it is the oriented surface
transfer dual to pullback on relative cohomology. Counting sheets
gives \(P_\nu T_\nu=4\).

Fill \(a\) in the base total and fill the positive transport of \(l\)
in the source total. Pull back the first filling and push forward
the second. A generic first filling avoids singular points of the
branch curve and meets its smooth part in finitely many transverse
points. Its full inverse image is an oriented branched surface
chain with no new boundary at those points. Move the other filling
off their preimages. Every intersection contributing to the ordered
pairing is now in the unramified locus. There the holomorphic map
preserves orientation, and the intersection projection formula is
the direct equality of signed local intersections on its sheets.
Because the map is over the value disk, pushforward also respects
positive nearby-page transport. This proves the first pairing
identity. Putting the pulled-back filling on the neighboring page
instead proves the second, with no transpose substitution.
The restriction to \(T_\nu H_\Q\) is \(4s_h\); nondegeneracy and
(6) give the asserted two-sided complement and projection.
\end{proof}

\section{Boundary vectors and the full local lattice}

The degree-four cover
\[
 (U,V)\longmapsto(U^4,V^4,UV),\qquad
 (U,V)\longmapsto(\omega U,\omega^{-1}V),\quad \omega^4=1,
\]
turns the cone function into
\(\widetilde P=UV(U^{4a}+V^{4b})\).
Let \(b_0\) be the complex-oriented boundary of \(V=0\), and let
\(b_1\) be the oppositely oriented boundary of \(U=0\), downstairs.
The two lens classes are both \(1\): their upstairs complex
orientations give respectively the two inverse deck generators,
and the second orientation was reversed.

\begin{lemma}\label{lem:boundary}
The form \(e\) is nondegenerate, its Seifert operator has finite
order, both \(b_j\) are fixed, and
\[
 \tau(b_j)=1,\quad e(b_j,b_j)=-n_j^2/4,\quad
 e(b_0,b_1)=-1/4,\quad
 4e(b_j,x)\equiv-\tau(x)\pmod4.                 \tag{7}
\]
The same boundary pairings hold in the reversed order.
\end{lemma}

\begin{proof}
The full transfer of an axis boundary is a single upstairs circle.
Different complex-oriented branch boundaries have pairing their
positive intersection number: fill them by their oriented complex
branches and use a small positive push. The sum of all boundary
classes on the page is zero. The first axis meets the other axis
once and \(U^{4a}+V^{4b}\) with total multiplicity \(4a\).
Its self-pairing is therefore \(-(1+4a)=-n_0^2\);
the second is \(-n_1^2\). Reversing the second boundary gives
mixed value \(-1\). Division by four proves the rational values.

The upstairs pairing of transfers is integral. If \(\tau(x)=0\),
represent \(x\) by a based loop with trivial deck monodromy.
Its transfer is the sum of four deck translates of a closed lift.
The transferred axis is deck invariant, so its four pairings with
these translates are equal. Thus \(4e(b_j,x)\) is divisible by four.
Evaluation at \(b_j\), together with \(n_j^2\equiv1\pmod4\),
identifies the induced functional on \(\Z/4\) as minus the identity.

Rational transfer identifies \(E_\Q\) with the deck invariants in
upstairs homology. Deck transformations preserve the upstairs
Seifert form. Averaging over them is self-adjoint in both orders,
so its image is nondegenerate and orthogonal to its kernel.
This proves nondegeneracy of \(e\).
Weighted rotation, with weights \(b/m,a/m\), has finite order
and commutes with the deck group. Its Seifert operator is this
rotation or its inverse, with the same convention as for \(h\).
It fixes each axis class. For any bilinear form \(s\) with
\(A=s^{-1}s^{\mathsf t}\), an \(A\)-fixed vector pairs
symmetrically with every vector. This gives the last assertion.
\end{proof}

The parameter arithmetic is exact:
\[
 \gcd(n_0,n_1)=1,\quad n_j\equiv1\pmod4,\quad
 a,b\equiv2\pmod5,\quad 5\nmid g,\quad 2^{r-1}\mid g,\quad g>5.
                                                               \tag{8}
\]
A common divisor of the \(n_j\) divides \(5\cdot2^r\), while neither
integer is even or divisible by five. The other assertions follow
by factoring \(n_j^2-1\).

Set \(v_j=b_j/n_j\). For any prime \(\ell\nmid n_j\), (7) gives
\[
 c_j(x)=-4e(v_j,x)\in\Z_\ell,\quad
 c_j(v_j)=1,\quad
 \tau(v_j)=1,\quad c_j(x)\equiv\tau(x)\pmod{4\Z_\ell}.             \tag{9}
\]
At odd primes the residue target is zero. At two, both last
equalities use \(n_j\equiv1\pmod4\), not merely that \(n_j\) is odd.
In particular
\[
 E_{\Z_\ell}=\Z_\ell v_j\perp U_{j,\ell},\qquad
 U_{j,\ell}=\ker c_j\subset\ker\tau.                             \tag{10}
\]
The kernel and coefficient are integral, although \(e\) itself
need not be integral over \(\Z_2\).

The localized version of (2) is exact, not an additional lattice
assumption. Write its congruence as the kernel of
\(D_\nu\oplus(H\otimes E)\to H/4H\), then take the image under
\((l,z)\mapsto(\pr_\nu l,z)\). Clearing the denominator four
makes the latter a map between finitely generated integer modules.
Smith normal form shows that tensoring these kernel and image
descriptions with \(\Z_\ell\) preserves them: \(\Z_\ell\) is
torsion-free over the principal ideal domain \(\Z\).
Thus (2) holds with exactly the same residue condition over
\(\Z_\ell\), including \(\ell=2\).

\begin{proposition}\label{prop:local}
For every \(\ell\nmid n_j\), and also over \(\Q\), the map
\[
 \alpha_{\nu,j}:(D_\nu,-s_{d_\nu})_{\Z_\ell}
       \perp(H\otimes U_{j,\ell},s_h\otimes e)
       \longrightarrow(\Lambda_\nu,s_\nu)_{\Z_\ell},\qquad
 (l,u)\longmapsto(\pr_\nu l,P_\nu l\otimes v_j+u)                 \tag{11}
\]
is an onto integral Seifert isometry.
\end{proposition}

\begin{proof}
Its image satisfies (2) by (9). Conversely decompose the second
coordinate of an arbitrary element of (2) as \(x\otimes v_j+u\).
The residue condition says \(x-P_\nu l\in4H_{\Z_\ell}\).
Replace \(l\) by
\[
 l'=l+T_\nu((x-P_\nu l)/4).
                                                               \tag{12}
\]
This is integral even at two: the numerator has already been
proved divisible by four. It has the same \(\pr_\nu l\) and
\(P_\nu l'=x\). Thus (11) is onto. Its rational kernel is zero
by \(P_\nu T_\nu=4\), proving injectivity.

Both cross terms vanish. From the two identities (6),
\[
 s_{d_\nu}(\pr_\nu l,\pr_\nu l')
   =s_{d_\nu}(l,l')-\tfrac14s_h(P_\nu l,P_\nu l').
 \]
The \(v_j\)-contribution is \(-\tfrac14s_h(P_\nu l,P_\nu l')\).
The leading minus sign on \(K_\nu\) in (2) makes their sum
exactly \(-s_{d_\nu}(l,l')\), with the given argument order.
\end{proof}

Use the abstract isometry \(I\) from Lemma~\ref{lem:plane} to define
\[
 J_j=\alpha_{1,j}(I\oplus1)\alpha_{0,j}^{-1}:F_0\longrightarrow F_1.
                                                               \tag{13}
\]
These are rational Seifert isometries, integral at every
\(\ell\nmid n_j\). Equation (12), not a rank or determinant count,
establishes that these are isometries of the full localized lattices.

\section{Discrepancy and a rational Clifford lift}

Write \(A_\nu=s_\nu^{-1}s_\nu^{\mathsf t}\) and
\(B_\nu=s_\nu+s_\nu^{\mathsf t}=s_\nu(1+A_\nu)\).
The two-sided decomposition in (6) makes \(K_\nu\) invariant.
Lemma~\ref{lem:plane} and the finite order of the two operators on
\(H\otimes E\) imply that \(A_\nu\) is semisimple at \(1\) and
the primitive fifth roots. Thus
\[
 F_{\nu,1}=\ker(A_\nu-1),\qquad
 F_{\nu,5}=\ker\Phi_5(A_\nu)
                                                               \tag{14}
\]
are rational two-sided orthogonal primary summands.
Indeed an isometry pairs generalized eigenvectors only at reciprocal
eigenvalues, as follows by applying its minimal polynomials in the
two arguments. Both sets of roots in (14) are inversion-stable.
On them \(1+A_\nu\) is invertible, so \(B_\nu\) is nondegenerate.

Let \(L=\Q(\zeta_5)\), \(L^+=\Q(\zeta_5+\zeta_5^{-1})\), acting
on \(F_{\nu,5}\) by \(\zeta_5=A_\nu\). Trace duality defines
the nondegenerate Hermitian form \(h_{\nu,5}\) by
\[
 \Tr_{L/\Q}(c\,h_{\nu,5}(x,y))=B_\nu(cx,y)\quad(c\in L).          \tag{15}
\]
The \(B_\nu\)-adjoint of \(c\) is \(\bar c\). This proves
sesquilinearity, Hermitian symmetry, and the equivalence of
\(L\)-linear \(B_\nu\)-isometries with isometries of (15).

\begin{lemma}\label{lem:lift}
For \(\Delta=J_0^{-1}J_1\), the action outside
\(F_{0,1}\oplus F_{0,5}\) is the identity.
Its first restriction has a rational Spin lift; its second lies
in \(\SU(F_{0,5},h_{0,5})(L^+)\).
\end{lemma}

\begin{proof}
Put \(w=v_0-v_1\). Equations (7) give
\[
 e(w,w)=-\tfrac12(1-1/(n_0n_1))\ne0.
 \]
The \(e\)-reflection
\(R(x)=x-2e(x,w)w/e(w,w)\) exchanges \(v_0,v_1\), preserves
\(e\), and commutes with its operator, since \(w\) is fixed.
Let \(\rho_\nu=1\) on \(K_\nu\) and \(1_H\otimes R\) on
\(H\otimes E\). In the coordinates (11), \(\rho_\nu\) is
the identity on the \(D_\nu\) coordinate and sends \(U_0\) to
\(U_1\) by the same \(R\) on both sides. Explicitly,
\(\rho_\nu\alpha_{\nu,0}(l,u)=\alpha_{\nu,1}(l,(1_H\otimes R)u)\).
The domain maps commute with \(I\oplus1\), so
\[
 J_1=\rho_1J_0\rho_0,\qquad
 \Delta=(J_0^{-1}\rho_1J_0)\rho_0.                              \tag{16}
\]
The identity follows from (11) on every \(l\) and \(u\);
it does not require \(P_1I=P_0\) or any transfer compatibility.

The minus space of \(\rho_\nu\) is
\(W_\nu=H_\Q\otimes\Q w\), with restricted form \(e(w,w)s_h\).
It has only the primary types \(1,5\).
Each involution fixes the two-sided orthogonal complement of
its minus space, giving the support assertion in (16).
The natural map \(x\otimes w\mapsto x\otimes w\) identifies the
minus spaces on the two sides, separately on each primary sector.
Composing it with \(J_0^{-1}\) identifies the minus spaces of
the two factors in (16), without identifying their embeddings.
Their dimensions are two over \(\Q\) on the first sector and
one over \(L\) on the fifth sector. They are nondegenerate,
because \(e(w,w)\ne0\) and the corresponding symmetric base
forms are nondegenerate.

The two determinants on the fifth sector are both \(-1\);
their product is one. By (15) this is exactly a special-unitary
element over \(L^+\).

There is also a direct rational lift on the first sector, so no
local spinor-norm computation is being substituted for a rational
lift. Put \(q(x)=B_0(x,x)\) on that sector. Choose a rational
orthogonal basis \(u_1,u_2\) of the minus space of \(\rho_0\),
and let \(u'_1,u'_2\) be its images in the minus space of
\(J_0^{-1}\rho_1J_0\) under the isometry just constructed.
Write \(q(u_i)=q(u'_i)=q_i\ne0\).
An involution with this minus space is the product of the
reflections in its two basis vectors. In the Clifford algebra
with \(v^2=q(v)\), the even element
\[
 z=\frac{u'_1u'_2u_1u_2}{q_1q_2}                               \tag{17}
\]
therefore acts as the first restriction of \(\Delta\).
Reversing the four factors shows directly that its Clifford norm
is \(q_1^2q_2^2/(q_1q_2)^2=1\). It normalizes the vector space,
as a product of four nonisotropic vectors, and thus belongs to
\(\Spin(q)(\Q)\). Scaling the quadratic form by a nonzero constant
does not change the required orthogonal action; our convention
uses \(q=B_0(x,x)\) throughout.
\end{proof}

\section{Real indefiniteness from actual holomorphic forms}

\begin{lemma}\label{lem:indices}
For each \(q\in\{1,4\}\) there are positive \(i,j\) such that
\[
 i\equiv j\pmod4,\qquad bi+aj=qm/5.                            \tag{18}
\]
They can be chosen with \(j\le4b/g\). Every such positive solution
has \(i<4a+1,\ j<4b+1\).
\end{lemma}

\begin{proof}
Set \(a=gA,b=gB\), where \(\gcd(A,B)=1\).
Let \(d=\gcd(A+B,4B)\); then \(d\mid4\).
Also \(5\mid m/g\), since \(a,b\equiv2\pmod5\) and \(5\nmid g\).
Since \(d\mid 4gAB+A+B=m/g\), coprimality with five gives
\(d\mid m/(5g)\).
Solve
\[
 (A+B)j+4Bt=qm/(5g),\qquad i=j+4t.
 \]
Choose \(1\le j\le4B/d\). The inequality
\[
 m/(5g)-4AB=\{4(g-5)AB+A+B\}/5>0
 \]
makes \(i=(qm/(5g)-Aj)/B\) positive. Finally
\[
 bi\le4m/5-a=b(4a+1)-m/5<b(4a+1),
 \]
and the symmetric calculation bounds \(j\).
\end{proof}

\begin{lemma}\label{lem:forms}
Each of the primitive \(\zeta_5,\zeta_5^{-1}\) eigenspaces of
\(E_\C\) has both signs for the Hermitian intersection form.
\end{lemma}

\begin{proof}
On the smooth affine curve
\[
 C:\quad UV(U^{4a}+V^{4b})=1
 \]
both \(U,V\) are nonzero. The weighted Euler identity proves
smoothness. For each of the two pairs in Lemma~\ref{lem:indices}
the differential
\[
 \omega_{i,j}
   =U^{i-1}V^{j-1}\frac{dU}{\partial_V\widetilde P}
   =-U^{i-1}V^{j-1}\frac{dV}{\partial_U\widetilde P}              \tag{19}
\]
is regular and not identically zero: the two expressions cover
the curve, and at a point where the first denominator is nonzero
the coordinate \(U\) is a local parameter. Its deck character is
\(\omega^{i-j}=1\). Weighted rotation has character
\(\exp(2\pi i(bi+aj)/m)=\zeta_5^q\).

We check every point at infinity on the smooth compactification
of \(C\), rather than assume that a regular affine differential
is a holomorphic cohomology class. Write \(u=\operatorname{ord}U\),
\(v=\operatorname{ord}V\). At least one coordinate has a pole:
otherwise the point would belong to the already smooth affine
curve. If \(4au<4bv\), the equation gives
\(u<0,\ v=-(4a+1)u\); the order of (19) is
\[
 (i-(4a+1)j)u-1\ge0.                                         \tag{20}
\]
This uses \(i<4a+1\) and \(j\ge1\), not a numerical estimate
at a sample parameter. The case \(4bv<4au\) is symmetric,
using the second expression in (19).
In the remaining case \(4au=4bv<0\), write
\((u,v)=-c(B,A)\), \(c\) a positive integer.
The leading cancellation in \(U^{4a}+V^{4b}\) does not cancel
\(\partial_V\widetilde P=U(U^{4a}+(4b+1)V^{4b})\);
its order is \(u+4bv\). Thus the order of (19) is
\[
 (i-1)u+(j-1-4b)v-1
   =c(1-q/5)m/g-1\ge0,                                      \tag{21}
\]
because \(m/(5g)\) is a positive integer. These three cases
exhaust all ends. Both differentials are therefore nonzero
holomorphic differentials on the compactification.

No Hodge decomposition theorem is needed for the signs.
For a nonzero holomorphic differential \(\omega=f(z)\,dz\),
\(i\omega\wedge\bar\omega=2|f(z)|^2\,dx\wedge dy\) has strictly
positive integral. Its period class is nonzero: if exact, Stokes'
theorem would make that integral zero. Cutting the surface along
a symplectic basis and applying Stokes to a primitive on the cut
polygon identifies this integral with the Hermitian cup pairing
of the periods. Complex conjugation reverses its sign.
Accordingly the holomorphic form of character \(\zeta_5\) and
the conjugate of the holomorphic form of character
\(\zeta_5^{-1}\) give opposite signs in the same character;
the other character is handled the same way. Passing by the
nondegenerate surface intersection pairing to homology changes
characters reciprocally and a common sign, neither of which
removes the two signs.

The compactification is a compact oriented surface with finitely
many punctures removed to obtain \(C\).
Inclusion induces an onto homology map preserving intersections;
its kernel consists of puncture-boundary classes.
The finite group generated by deck transformations and weighted
rotation acts on this map. Average any linear section over that
group, so the lifted classes have the prescribed characters and
unchanged intersection numbers. Rational transfer to the deck
invariants identifies these with classes in \(E_\C\), with a
positive covering-degree factor on intersection.
To identify the affine homology with the local weighted page,
positive weighted scaling takes a small nonzero fiber to \(C\);
as the value decreases its cut exhausts \(C\).
Outside a compact core every end is an annulus, so such a cut
induces precisely the same first homology and intersections.
This also makes the weighted rotation action the local one.
Thus the claimed two signs occur in the actual cone-page group.
\end{proof}

\begin{proposition}\label{prop:real}
\(B_\nu|F_{\nu,1}\) is indefinite and has dimension at least six.
The Hermitian space in (15) has \(L\)-dimension at least four
and is indefinite at the real place of \(L^+\) corresponding to
\(\{\zeta_5,\zeta_5^{-1}\}\).
\end{proposition}

\begin{proof}
Let \(\lambda\) be \(\zeta_5\) or its inverse, in the chosen
Seifert-operator convention. Choose a nonzero base vector
\(x\in H_\C\) with \(A_hx=\lambda x\), and a cone vector
\(y\) with \(A_ey=\lambda^{-1}y\). Put
\[
 I_H=i(s_h-s_h^{\mathsf t})(x,\bar x),\qquad
 I_E=i(e-e^{\mathsf t})(y,\bar y).
 \]
The base eigenspace is one-dimensional and pairs nondegenerately
with its inverse, so \(I_H\) is real and nonzero.
Lemma~\ref{lem:forms} supplies both signs of the real number \(I_E\).
Using \(s^{\mathsf t}=sA\) in both factors gives exactly
\[
 B_\nu(x\otimes y+\bar x\otimes\bar y,\
             x\otimes y+\bar x\otimes\bar y)
   =-\frac{4I_HI_E}{(1-\lambda)(1-\lambda^{-1})}.                \tag{22}
\]
The denominator is \(|1-\lambda|^2>0\). These are real fixed
vectors, so both quadratic signs occur.
The two signs in one complex cone eigenspace force its dimension
to be at least two; tensors with the base line and their conjugates
give at least four real dimensions. The independent fixed tensor
subspace \(H^{(1)}\otimes\Q b_0\) gives two more.

Choose a rational fixed base vector \(x_0\) with
\(s_h(x_0,x_0)\ne0\), possible because its symmetric restriction
on the two-dimensional fixed subspace is nondegenerate.
On a cone eigenspace with eigenvalue \(\lambda\),
\[
 (e+e^{\mathsf t})(y,\bar y)=
  \frac{1+\lambda^{-1}}{i(1-\lambda^{-1})}
          i(e-e^{\mathsf t})(y,\bar y).                         \tag{23}
\]
The ratio is real and nonzero, since \(\lambda\ne1,-1\).
Tensoring with \(x_0\) thus gives both signs for \(B_\nu\)
in the desired primitive component. Trace duality (15) identifies
this with the Hermitian form at that real place, up to its positive
trace normalization. It is indefinite. Moreover the whole
two-dimensional fixed base space tensor the cone eigenspace of
dimension at least two embeds in this primitive component.
Its complex dimension, and hence the \(L\)-dimension in (15),
is at least four.
\end{proof}

\section{The global correction preserves the full lattice}

We state exactly the new external input. Rapinchuk's versioned
primary text \cite[Theorem 2.3]{rapinchuk} says: for a connected
absolutely almost simple group over a global field and a finite
set \(S\) of places, strong approximation away from \(S\) holds
if and only if the group is simply connected and
\(\prod_{v\in S}G(K_v)\) is noncompact. We use only number fields
and omit all archimedean places.

Consider
\[
 G_1=\Spin(F_{0,1},B_0)\quad\hbox{over }\Q,\qquad
 G_5=\SU(F_{0,5},h_{0,5})\quad\hbox{over }L^+.
 \]
The classical-group facts are primary-text inspected, with their
precise scope recorded in the accompanying excerpts:
\begin{itemize}
\item Conrad's nondegenerate quadratic-space results give the
simply connected central cover \(\Spin(q)\to\mathrm{SO}(q)\)
with kernel \(\mu_2\). Over the algebraic closure, its root
system is \(B_m\) for dimension \(2m+1\), or \(D_m\) for
dimension \(2m\) \cite[C.3.10, C.4.1, C.4.10]{conrad}.
In dimension at least five these systems are irreducible;
dimension six is \(D_3=A_3\), not the excluded \(D_2\) case
\cite[C.6.6]{conrad}. Thus \(G_1\) is connected, simply
connected, and absolutely almost simple.
\item For a nondegenerate Hermitian space take
\(\operatorname{End}_L(V)\) with its adjoint involution.
After algebraic closure of \(L^+\) it is a product of two matrix
algebras with exchange-transpose involution
\cite[24.40(b)]{milne}.
The unitary equations give \(Y=(X^{\mathsf t})^{-1}\);
the determinant-one condition gives \(X\in\mathrm{SL}_d\).
The actual bound \(d\ge4\) also meets the \(n>2\) heading of
the subsection containing Milne's equation (158); no rank-two
outer-form classification is needed. The displayed equations
and the assertion for \(\mathrm{SL}_d\) in Summary 21.96 give connectedness,
simple connectedness and absolute almost simplicity of \(G_5\).
\end{itemize}
Noncompactness is not a citation inference. An indefinite real
quadratic space contains orthogonal \(u,v\) with \(q(u)=1\),
\(q(v)=-1\); then \(uv\) squares to \(1\).
The Clifford elements
\(\cosh(t/2)+\sinh(t/2)uv\) have norm one and act by
unbounded boosts. This proves \(G_1(\R)\) noncompact.
An indefinite complex Hermitian space contains a
\(\operatorname{diag}(1,-1)\) plane. The matrices
\(\left(\begin{smallmatrix}\cosh t&\sinh t\\
\sinh t&\cosh t\end{smallmatrix}\right)\), extended by the identity,
belong to its special unitary group and are unbounded.
Proposition~\ref{prop:real} supplies this plane at the required
real place of \(L^+\).

Apply Rapinchuk's theorem separately over \(\Q\) and \(L^+\).
Restriction of scalars and a finite product give density of
\[
 G(\Q)\text{ in }G(\mathbb A_f),\qquad
 G=G_1\times\operatorname{Res}_{L^+/\Q}G_5.                     \tag{24}
\]
This does not incorrectly call \(G\) absolutely almost simple
over \(\Q\). The adelic identification for the restriction of
scalars follows place by place from
\(L^+\otimes\Q_\ell=\prod_{v\mid\ell}L^+_v\).

Let \(\rho\) be the action of \(G\) on the two primary spaces,
extended as the identity on the two-sided primary complement.
It preserves \(B_0\) and commutes with \(A_0\); since
\(s_0=B_0(1+A_0)^{-1}\) on these sectors, it preserves \(s_0\).
Lemma~\ref{lem:lift} gives
\(\widetilde\Delta\in G(\Q)\) mapping to \(\Delta\).
At each rational prime define
\[
 a_\ell=\begin{cases}\widetilde\Delta&\ell\mid n_0,\\1&\ell\nmid n_0.
 \end{cases}
 \qquad
 C_\ell=\{\gamma\in G(\Q_\ell):
       \rho(\gamma)\Lambda_{0,\Z_\ell}=\Lambda_{0,\Z_\ell}\}.
                                                               \tag{25}
\]
Here \(C_\ell\) is the stabilizer of the \emph{whole} lattice (2).
No integral primary decomposition is assumed.

We give the required adelic neighborhood explicitly enough to
exclude an implicit integrality gap. Choose a faithful rational
matrix representation \(G\subset\mathrm{GL}_N\) and a basis
of \(\Lambda_0\). The entries of \(\rho\) and \(\rho^{-1}\)
are regular functions in the group coordinates. Clearing their
finitely many rational coefficients, including those of a
presentation of \(G\), gives an integer \(D\) such that
\[
 K_\ell^0=G(\Q_\ell)\cap\mathrm{GL}_N(\Z_\ell)\subset C_\ell
 \quad(\ell\nmid D).
 \]
These are the standard compact opens for this faithful
representation. At each of the finitely many remaining primes,
\(C_\ell\) is open by continuity of \(\rho\) and contains the
identity. A sufficiently small matrix congruence subgroup
\(K_\ell\subset K_\ell^0\) therefore lies in \(C_\ell\).
Use \(K_\ell=K_\ell^0\) at all other primes. Then
\(\mathcal K=\prod_\ell K_\ell\) is compact open in \(G(\mathbb A_f)\).
This argument only needs compact opens \emph{inside} \(C_\ell\);
it does not assume a faithful Spin action on \(F_{0,1}\).

The finite-support tuple \(\mathbf a=(a_\ell)_\ell\) is adelic,
and \(\mathbf a\mathcal K\) is a nonempty open set. By (24)
choose \(g_{\rm rat}\in G(\Q)\cap\mathbf a\mathcal K\).
For \(\ell\nmid n_0\), \(J_0\) is integral by
Proposition~\ref{prop:local}. For \(\ell\mid n_0\),
coprimality in (8) makes \(J_1\) integral and
\(J_0\rho(a_\ell)=J_0\Delta=J_1\).
Writing \(g_{{\rm rat},\ell}=a_\ell k_\ell\) with \(k_\ell\in K_\ell\),
at every prime we have
\[
 J_0\rho(g_{\rm rat})\Lambda_{0,\Z_\ell}
     =J_0\rho(a_\ell)\rho(k_\ell)\Lambda_{0,\Z_\ell}
     =\Lambda_{1,\Z_\ell}.                                    \tag{26}
\]
In particular the correction at two lies in the full-lattice
stabilizer, preserving the mod-four residue condition.
Finally a rational vector in every localization of a full
integer lattice has integral coordinates in an integer basis:
each rational coordinate has no prime in its denominator.
Intersecting (26) at all primes proves
\[
                    J_0\rho(g_{\rm rat})\Lambda_0=\Lambda_1.
 \]
This is an onto integral Seifert isometry and proves
Theorem~\ref{thm:result}.

\section{What the exact replay does, and does not, certify}

Run \texttt{python3 research/m45\_arithmetic/verify.py}.
The replay checks the source-bound polynomial and branch identities,
the annihilating-polynomial implication in (5), the boundary Gram
matrix and reflection, the local pairing cancellation, Clifford
norm cancellation in (17), the eigenform order identities, and
the sign identities (22)--(23). It also constructs positive indices
for bounded admissible parameter examples and checks the first
nontrivial pair exactly. Those examples are regression checks,
not the proof of the universally quantified arithmetic lemma.
The largest actual matrix in these checks is two by two.

The Jacobian adapter was inspected before use. Its one bounded
operation was
\begin{center}
\small\texttt{matrix.symbolic.characteristic\_polynomial.compute}
\end{center}
on the actual two-boundary Gram matrix
\(-\tfrac14\left(\begin{smallmatrix}n_0^2&1\\1&n_1^2\end{smallmatrix}\right)\).
The exact payload, typed output, server/protocol metadata and exit
status are retained in the dedicated \texttt{jacobian\_boundary}
JSON files. The local replay independently recomputes that result;
no timeout or unknown result is counted as mathematical evidence.

The following are \emph{prose proof obligations}, not conclusions
of a finite matrix experiment: the simultaneous plane resolution
and ordered-form transport; the branched two-sided intersection
projection formula; the two-annulus monodromy construction;
identification of the actual boundary classes and transfer;
exhaustion of all compactification ends and nonzero period classes;
the classical-group facts from the inspected primary sources;
the applicability of strong approximation and the full-lattice
adelic intersection argument. Their proofs and precise inputs
are given above so that an independent Reviewer can assess them.

Relative to the strongest prior result, this supplies the missing
comparison \emph{between} the two reviewed integral tube lattices.
It does not replace the outstanding classification application
by equality of Milnor numbers, by local equivalence, or by this
compact computational certificate.

\begin{thebibliography}{9}
\bibitem{rapinchuk}
A. S. Rapinchuk, \emph{On strong approximation for algebraic groups},
arXiv:1207.4425v1, Theorem 2.3,
\url{https://arxiv.org/html/1207.4425v1}.
\bibitem{conrad}
B. Conrad, \emph{Reductive group schemes}, 2014 author manuscript,
Appendix C, especially C.3.10, C.4.1, C.4.10 and C.6.6,
\url{https://math.stanford.edu/~conrad/papers/luminysga3.pdf}.
\bibitem{milne}
J. S. Milne, \emph{Algebraic Groups}, corrected 2022 CUP reprint,
Proposition 24.40(b), equation (158), Summary 21.96, and the
Hermitian adjoint construction on printed p.~538,
\url{https://www.jmilne.org/math/Books/iAG2022.pdf}.
\end{thebibliography}
\end{document}
